\documentclass[11pt]{article}
\usepackage[no-math]{fontspec}
\usepackage[french,english]{babel}
\usepackage{amsmath}
\usepackage{amssymb}
\usepackage{amsthm}
\usepackage{mathrsfs}
\usepackage{tikz-cd}
\usepackage{enumitem}

\newcommand{\printedpage}[1]{\marginpar{\footnotesize p.~#1}}
\newcommand{\msp}[1]{\marginpar{\footnotesize [#1]}}
\newcommand{\ul}[1]{\underline{\mathrm{#1}}}
\newcommand{\xline}[1]{\mathrel{\overset{#1}{\rule[0.6ex]{2em}{0.4pt}}}}

\begin{document}

\printedpage{4}
% head: A. Grothendieck : Dérivateurs, Ch. VI, Édition des Universités de Montpellier II et Paris VII
and not only transitivity up to canonical isomorphism. Thus one can also write
\[
\mathrm{HOT}_W(S) \xrightarrow{\ \sim\ } (\ul{Fibsc}\,S)W_S^{-1}, \tag*{$(*)$}
\]
where this time $W_S$ is interpreted as
\[
W_S \subseteq \mathrm{Fl}(\ul{Fibsc}\,S),
\]
namely
\[
W_S = \{u \in \mathrm{Fl}(\ul{Fibsc}\,S) \mid u_s \in W \text{ for all } s \in \mathrm{Ob}\,S\}
\]
(strictly cartesian functors, i.e. commuting with the splittings, which induce $W$-equivalences between the fiber categories).

The arrow in $(*)$ is indeed an isomorphism of categories, not just an equivalence. But we are going to define other categories, more convenient in certain respects, which will be \textit{equivalent} (we hope), but certainly not isomorphic, to $\mathrm{HOT}_W(S)$. (We shall deal later with the variance of $\mathrm{HOT}_W(S)$ in $S$.)

The categories $\ul{Hom}(S^o,\mathrm{Cat})$, $\ul{Fibsc}\,S$, as well as the variants we are now going to introduce, are in fact 2-categories. But to describe $\mathrm{HOT}_W(S)$, only the 1-category structures are used. One has the impression that this is quite % [sic] seule les structures -> seules les structures, accord du pluriel
\msp{page 5}
crude, that one loses something in doing so. But what is lost will be recovered, and superabundantly, in the structure of the internal \underline{Hom} in the $\mathrm{HOT}(S)$, if one ``remembers'' that these internal \underline{Hom}, giving rise to ``relative homotopy types'' $\ul{Hom}_{\mathrm{HOT}_W}(S;\xi,\eta)$ (if $\xi$, $\eta$ belong to them), can be interpreted in terms of $\infty$-stacks in groupoids over $S$, with all the richness of structure that belongs to them. But this is another story\dots

In fact, I see a diagram of 13 categories of ``categories over $S$\,'', to be examined now % [sic] categories -> catégories, accent manquant
\[
\begin{tikzcd}[column sep=small]
 & \ul{Fibsc}_0S \arrow[dl, hook] \arrow[d, "\mathrm{fid}"] & & \ul{Cofibsc}_0S \arrow[d] \arrow[dr, hook] & \\
\ul{Fibsc}\,S \arrow[d, "\mathrm{fid}"] & \ul{Fib}_0S \arrow[dl, hook] & & \ul{Cofib}_0S \arrow[dr, hook] & \ul{Cofibsc}\,S \arrow[d, "\mathrm{fid}"] \\
\ul{Fib}\,S \arrow[dr, hook] & & \ul{Parf}\,S \arrow[dl, hook] \arrow[d, hook] \arrow[dr, hook] & & \ul{Cofib}\,S \arrow[dl, hook] \\
 & \ul{Liss}\,S \arrow[dr, hook] & \ul{Fibhot}\,S \arrow[d, hook] & \ul{Prop}\,S \arrow[dl, hook] & \\
 & & \mathrm{Cat}/S & &
\end{tikzcd}
\]
% foot: Version 9 novembre 2001

\printedpage{5}
% head: A. Grothendieck : Dérivateurs, Ch. VI, Édition des Universités de Montpellier II et Paris VII
where \underline{Fib}, \underline{Cofib} denote the Cat-fibered, Cat-cofibered categories over $S$, \underline{Liss} and \underline{Prop} the categories
\msp{page 6}
smooth resp. proper over $S$, \underline{Parf} those that are both smooth and proper (perfect fibrations), finally \underline{Fibhot} the Hot-fibered categories over $S$. All the arrows are strictly fully faithful functors, except for the two (rather four) upper vertical arrows, which are faithful and no more. I have just added another 4 of them, categories: $\ul{Fibsc}_0$, $\ul{Fib}_0$, $\ul{Cofibsc}_0$, $\ul{Cofib}_0$ are the full subcategories of the corresponding categories without subscript, consisting of the fibered resp. cofibered categories $X$ over $S$, for which the base change resp. cobase change functors
\[
u^* : X_t \longrightarrow X_s, \qquad\quad \text{resp.} \qquad\quad u_* : X_s \longrightarrow X_t,
\]
for $s \xrightarrow{u} t$ in $S$, are Hot-equivalences. In terms of functors
\[
S^o \xrightarrow{F} \mathrm{Cat} \qquad\quad \text{resp.} \qquad\quad S \xrightarrow{F'} \mathrm{Cat},
\]
this means that this functor transforms arbitrary arrows in $S$ into elements of $W$ in Cat, i.e. that we have a commutative diagram
\[
\begin{tikzcd}[row sep=large]
S^o \arrow[r, "F"] \arrow[d, "\mathrm{can}"'] & \mathrm{Cat} \arrow[d, "\mathrm{can}"] \\
\Pi(S^o) \simeq (\Pi S)^o = S^o(\mathrm{Fl}\,S^o)^{-1} \arrow[r, "\bar F"'] & \mathrm{Hot}_W = (\mathrm{Cat})W^{-1}
\end{tikzcd}
\]
(and likewise for $F'$), i.e. that we have a ``local
\msp{page 7}
system'' $\bar F$ on $S$, with values in $\mathrm{Hot}_W$, defined using $F$. I recall that one also finds such a local system in the case of a Hot-fibered category over $S$ (and even a local system with values in Hot itself, instead of $\mathrm{Hot}_W$ - we tacitly use here that $W_\infty \subseteq W$, i.e. \fbox{L3 quater} (\footnote[5]{We shall assume it true throughout what follows.})). And the same holds a fortiori for the objects of $\ul{Parf}\,S$. For $\ul{Liss}\,S$ resp. $\ul{Prop}\,S$, one simply finds a functor
\[
S^o \xrightarrow{\bar F} \mathrm{Hot} \qquad\quad \text{resp.} \qquad\quad S \xrightarrow{\bar F'} \mathrm{Hot},
\]
whence, by composing with $\mathrm{Hot} \longrightarrow \mathrm{Hot}_W$, functors with values in $\mathrm{Hot}_W$.

In each of these 14 \texttt{[13?]} categories, we introduce a set of arrows, which will serve to localize, and we shall always denote it by the same notation $W_S$. In all these categories except only $\mathrm{Cat}/S$, $W_S$ denotes $W$-equivalence ``fiber by fiber''. In $\mathrm{Cat}/S$ on the other hand, it will be just $W$-equivalence. All the arrows of the diagram are compatible with the localizing sets of arrows.
% foot: Version 9 novembre 2001

\printedpage{6}
% head: A. Grothendieck : Dérivateurs, Ch. VI, Édition des Universités de Montpellier II et Paris VII
\msp{page 8}
\[
\begin{tikzcd}[column sep=small, row sep=large]
 & \ul{Fibsc}_0S \arrow[dl, hook, "\alpha'"'] \arrow[d] & & \ul{Cofibsc}_0S \arrow[d] \arrow[dr, hook, "\beta'"] & \\
\ul{Fibsc}S \arrow[d] & \ul{Fib}_0S \arrow[dl, hook, "\alpha"'] \arrow[ddr, dotted, bend right=20] & & \ul{Cofib}_0S \arrow[dr, hook, "\beta"] \arrow[ddl, dotted, bend left=20] & \ul{Cofibsc}\,S \arrow[d] \\
\ul{Fib}S \arrow[dr, hook] & & \ul{Parf}\,S \arrow[dl, hook] \arrow[d, hook] \arrow[dr, hook] & & \ul{Cofib}\,S \arrow[dl, hook] \\
 & \ul{Liss}\,S \arrow[dr, hook] & \ul{Fibhot}\,S \arrow[d, hook] \arrow[uuul, dotted, bend left=10, "\Phi_0"] \arrow[uuur, dotted, bend right=10, "\Psi_0"'] & \ul{Prop}\,S \arrow[dl, hook] & \\
 & & \mathrm{Cat}/S \arrow[uuull, dotted, bend left=60, "\Phi"] \arrow[uuurr, dotted, bend right=60, "\Psi"'] & &
\end{tikzcd}
\]
% ?? uncertain: starting point of the dotted arrows labelled Phi_0 and Psi_0 (the dotted lines run alongside Parf S down to Fibhot S) -- Fibhot S or Parf S

\texttt{[In this diagram, there are also} $\ul{Liss}_0$ \texttt{and} $\ul{Prop}_0$ \texttt{but without connection to the other objects.]}

I am almost sure of the

\textbf{Theorem.} (\footnote[6]{This is totally botched, cf. comments page 32 ff., solution of the perplexities pages 41 ff. See especially the recapitulatory Theorem-Scholie p. 70-81.}) \textit{All the preceding functors induce equivalences between the localized categories.} % [sic] Theorème -> Théorème, accent manquant

As there are 16 functors, it looks like an exercise in patience. The general principle is that one manages to define functors in the opposite direction, compatible with the $W_S$, in such a way that the composites are linked to the identity functors by arrows which are in $W_S$, so that by passage to the localizations these functors induce (equivalences) quasi-inverses of the functors deduced from those of the diagram. All these functors in the opposite direction arise essentially by composition, from three very familiar functors. Two of these are indicated by dotted lines on the diagram, as $\Phi$, $\Psi$. % [sic] exercise -> exercice, orthographe française ; % [sic] reliées -> reliés, accord avec « les composés »

\msp{page 9}
The third is a canonical functor
\[
C : \mathrm{Cat}/S \longrightarrow \ul{Parf}\,S,
\]
which we are also going to make explicit.
% foot: Version 9 novembre 2001

\printedpage{7}
% head: A. Grothendieck : Dérivateurs, Ch. VI, Édition des Universités de Montpellier II et Paris VII
The functor $\Phi$ (resp. $\Psi$) applied to $\ul{Fib}_0S$ (resp. $\ul{Cofib}_0S$), induces a functor
\[
\ul{Fib}_0S \longrightarrow \ul{Fibsc}_0S \qquad\quad \text{resp.} \qquad\quad \ul{Cofib}_0S \longrightarrow \ul{Cofibsc}_0S \tag*{$(*)$}
\]
(which will have to be checked). This provides all the ``inverse'' arrows we need, with the sole exception of $\alpha$, $\beta$, $\alpha'$, $\beta'$. But if we have proved, by the indicated method, that all the arrows of the diagram induce equivalences on the localized categories, to prove that the same is still true for $\alpha$, $\beta$, $\alpha'$, $\beta'$, it suffices to prove it for the two functors
\[
\ul{Fibsc}_0S \longrightarrow \ul{Fib}\,S, \qquad\qquad \ul{Cofibsc}_0S \longrightarrow \ul{Cofib}\,S,
\]
so it remains to define the functors in the opposite direction
\[
\ul{Fib}\,S \longrightarrow \ul{Fibsc}_0S, \qquad\qquad \ul{Cofib}\,S \longrightarrow \ul{Cofibsc}_0S,
\]
and as the sources are full subcategories of $\mathrm{Cat}/S$, it will suffice to define
\[
\Phi_0 : \mathrm{Cat}/S \longrightarrow \ul{Fibsc}_0S, \qquad\qquad \Psi_0 : \mathrm{Cat}/S \longrightarrow \ul{Cofibsc}_0S.
\]
For this, we shall set
\[
\Phi_0 = \Phi r C, \qquad\qquad \Psi_0 = \Psi r C, \tag*{$(**)$}
\]
\msp{page 10}
where $C : \mathrm{Cat}/S \longrightarrow \ul{Parf}\,S$, $\Phi$ and $\Psi$ are the 3 ``retrograde'' functors inverted above, and $r : \ul{Parf}(S) \hookrightarrow \mathrm{Cat}/S$ is the inclusion. We must therefore check that $\Phi_0$, $\Psi_0$ defined by $(**)$, which a priori have as target categories $\ul{Fibsc}\,S$ and $\ul{Cofibsc}\,S$, take their value in the subcategories with subscript 0.

I am going to give a single equivalence statement, which will cover (almost) all 16 cases at once:

\textbf{Corollary.} \textit{Let us consider the diagram}
\[
\begin{tikzcd}[column sep=small, row sep=small]
 & \ul{Parf}\,S \arrow[dd, "r"'] & \\
\ul{Fibsc}_0S \arrow[dr, "p"'] & & \ul{Cofibsc}_0S \arrow[dl, "q"] \\
 & \mathrm{Cat}/S &
\end{tikzcd}
\]
\textit{and let $\mathcal{F}$ be a full subcategory of $\mathrm{Cat}/S$, such that one of the functors $p$, $q$, $r$ factors through $\mathcal{F}$. Let $W_{\mathcal{F}}$ be induced on $\mathcal{F}$ by $W_S$ in $\mathrm{Cat}/S$. Then $\mathcal{F} \hookrightarrow \mathrm{Cat}/S$ induces an equivalence on the localized categories.}

\msp{page 11}
\textsc{Proof}. I say that it suffices to prove, for the three functors $p$, $q$, $r$, and for the functors in the opposite direction
% foot: Version 9 novembre 2001
\printedpage{8}
% head: A. Grothendieck : Dérivateurs, Ch. VI, Édition des Universités de Montpellier II et Paris VII
\[
\begin{tikzcd}[column sep=small, row sep=small]
 & \ul{Parf}\,S & \\
\ul{Fibsc}_0S & & \ul{Cofibsc}_0S \\
 & \mathrm{Cat}/S \arrow[uu, "C"] \arrow[ul, "\Phi_0"] \arrow[ur, "\Psi_0"'] &
\end{tikzcd}
\]
that the six composites
\[
p\Phi_0, \Phi_0p, q\Psi_0, \Psi_0q, rC, Cr
\]
are linked to the identity functors (in the relevant category) by an arrow which is in the corresponding localizer. For each of these three pairs of functors $(p,\Phi_0)$, $(q,\Psi_0)$, $(r,C)$, this will already prove that on passing to the localizations, they induce equivalences quasi-inverse to one another. But we can do better with the help of the following formal lemma.

\textbf{Lemma.} \textit{Let $(C,W)$, $(C',W')$ be two categories with localizers $W \subseteq \mathrm{Fl}\,C$, $W' \subseteq \mathrm{Fl}\,C'$. Let}
\[
p : C \longrightarrow C'
\]
\textit{be a functor compatible with the localizers, and suppose there exists a functor $\varphi : C' \longrightarrow C$ likewise compatible, and such that we have arrows (of unspecified direction)}
\[
\varphi p \xline{u} id_C, \qquad\qquad p\varphi \xline{v} id_{C'},
\]
\textit{which are in $W$ resp. $W'$ ``argument by argument''.}

\msp{page 12}
\textit{Then $p$ and $\varphi$ induce equivalences quasi-inverse to one another on passing to the localized categories (they have no merit). Moreover, for any full subcategory $C'_0$ of $C'$, such that $p$ factors (as $p_0 : C \longrightarrow C'_0$) through $C'_0$, if we denote by $W'_0$ the localizer induced in $C'_0$ by $C'$, then $p_0 : C \longrightarrow C'_0$ and $\varphi_0 = \varphi|C'_0 : C' \longrightarrow C$} \texttt{[rather} $C'_0 \longrightarrow C$\texttt{]} \textit{also induce equivalences quasi-inverse to one another on the localizations.} % [sic] merite -> mérite, accent manquant

Indeed, $\varphi_0p_0 = (\varphi i)p_0 = \varphi(ip_0) = \varphi p \xline{u} \mathrm{id}_C$ (here $i : C'_0 \hookrightarrow C'$ denotes the inclusion), and $F \overset{\text{def}}{=} p_0\varphi_0 = p_0\varphi i$ satisfies
\[
iF = \underbrace{ip_0}_{p}\varphi i = (p\varphi)i.
\]
To define $F \xline{v_0} \mathrm{id}_{C'_0}$, which is to be in $W'_0$ argument by argument,
\msp{page 13}
it suffices to define
\[
\underbrace{iF}_{=\,(p\varphi)i} \xline{\bar v_0} i,
\]
which is to be in $W'$ argument by argument, since $W'_0$ induced by $W'$ and since $F \longmapsto iF$
\[
\ul{Hom}(C'_0,C'_0) \longrightarrow \ul{Hom}(C'_0,C')
\]
% foot: Version 9 novembre 2001

\end{document}
